% Unit 075 terjemahan Bahasa Indonesia dari chapter6.tex baris 740--868.
% Sumber: Methods of Algebra, Volume 2, commit
% 9a5803ff77dd3257484cb177f851a73770a59dd3 (CC BY 4.0).
% stable-unit-id: o014.aljabr2.chapter6.induced-modules
% Cakupan: Bagian 6.4 lengkap.

% segment-id: o014.li2.chapter6.induced-modules.h001
\section{Modul terinduksi}\label{sec:induced-module}
\hypertarget{o014.aljabr2.chapter6.induced-modules}{ }

% segment-id: o014.li2.chapter6.induced-modules.p001
Untuk homomorfisme grup $\varphi:H\to G$,
Proposisi~\ref{prop:pullback-two-adjoint} memberikan adjoin kanan
$\Hom_H(\Bbbk[G],\cdot)$ dan adjoin kiri
$\Bbbk[G]\dotimes{\Bbbk[H]}(\cdot)$ bagi $\varphi^*$. Bagian ini berfokus
pada kasus umum ketika $H\subset G$ merupakan subgrup dan $\varphi$ adalah
homomorfisme penyertaan. Kita mulai dengan memperkenalkan notasi.

% segment-id: o014.li2.chapter6.induced-modules.d001
\begin{definition}[Modul terinduksi]\label{def:induced-module}
	\index{modul-$G$!terinduksi (induced)}
	\index[sym1]{Ind-ind@$\Ind^G_H$, $\iInd^G_H$}
	Misalkan $H$ subgrup dari $G$. Untuk modul-$H$ $N$, definisikan
	modul-$G$
	\[ \Ind^G_H(N):=\Hom_H(\Bbbk[G],N),\qquad
	\iInd^G_H(N):=\Bbbk[G]\dotimes{\Bbbk[H]}N. \]
	Ketika $N$ bervariasi, keduanya memberikan funktor
	$H\dcate{Mod}\to G\dcate{Mod}$.
\end{definition}

% segment-id: o014.li2.chapter6.induced-modules.p002
Kedua operasi tersebut secara umum dapat disebut \emph{induksi} dari $H$ ke
$G$. Kadang-kadang $\iInd^G_H$ disebut induksi dengan dukungan hingga; istilah
ini akan dijelaskan oleh Lema~\ref{prop:Ind-as-mappings} di bawah.

% segment-id: o014.li2.chapter6.induced-modules.t001
\begin{proposition}\label{prop:Ind-preservation}
	Funktor $\Res^G_H:G\dcate{Mod}\to H\dcate{Mod}$ serta
	$\Ind^G_H,\iInd^G_H:H\dcate{Mod}\to G\dcate{Mod}$ semuanya eksak.
	Funktor $\Ind^G_H$ mempertahankan objek injektif,
	$\iInd^G_H$ mempertahankan objek projektif, sedangkan $\Res^G_H$
	mempertahankan keduanya.
\end{proposition}
\begin{proof}
	Kita sudah mengetahui bahwa $\Res^G_H$ eksak. Dengan memilih wakil koset,
	terlihat bahwa $\Bbbk[G]$ merupakan modul-$\Bbbk[H]$ kiri bebas, sehingga
	$\Hom_{\Bbbk[H]}(\Bbbk[G],\cdot)$ eksak. Dengan cara serupa,
	$\Bbbk[G]$ merupakan modul-$\Bbbk[H]$ kanan bebas, sehingga
	$\Bbbk[G]\dotimes{\Bbbk[H]}(\cdot)$ eksak. Berdasarkan relasi adjungsi,
	pernyataan mengenai pemertahanan objek injektif dan projektif merupakan
	penerapan langsung Proposisi~\ref{prop:adjoint-injective-projective}.
\end{proof}

% segment-id: o014.li2.chapter6.induced-modules.p003
Proposisi~\ref{prop:Ind-preservation} memberikan cara membangun resolusi
injektif (atau projektif) dalam $G\dcate{Mod}$: ambil resolusi injektif (atau
projektif) dari modul-$G$ $M$ yang diberikan dalam $\Bbbk\dcate{Mod}$, lalu
induksikan resolusi itu ke $G$ dengan $\Ind^G_{\{1\}}$ (atau
$\iInd^G_{\{1\}}$).

% segment-id: o014.li2.chapter6.induced-modules.p004
Hasil berikut lazim disebut \emph{Lema Shapiro}.

% segment-id: o014.li2.chapter6.induced-modules.t002
\begin{theorem}[B.\ Eckmann, D.\ K.\ Faddeev, A.\ Shapiro]
\label{prop:Shapiro}
	\index{Lema Shapiro (Shapiro's Lemma)}
	Misalkan $H$ subgrup dari $G$. Untuk setiap modul-$H$ $N$ dan
	$n\in\Z$, terdapat isomorfisme kanonik
	\[ \Hm^n\left(G,\Ind^G_H(N)\right)\simeq\Hm^n(H,N),\qquad
	\Hm_n\left(G,\iInd^G_H(N)\right)\simeq\Hm_n(H,N). \]
	Kedua sisi setiap isomorfisme, sebagai funktor dalam peubah $N$,
	mempunyai barisan eksak panjang, dan isomorfisme tersebut kompatibel dengan
	barisan-barisan eksak panjang itu.
\end{theorem}
\begin{proof}
	Kita berikan dua metode. Untuk kohomologi, metode pertama meninjau funktor
	\[ A^n:=\Hm^n\left(G,\Ind^G_H(\cdot)\right). \]
	Karena $\Ind^G_H$ eksak, keluarga funktor ini mempunyai barisan eksak
	panjang yang berasal dari $\Hm^n(G,\cdot)$, sehingga menjadi funktor-$\delta$
	kohomologis dalam arti Definisi~\ref{def:delta-functor}. Untuk $n=0$,
	Korolari~\ref{prop:Ind-invariant} memberikan isomorfisme kanonik
	\[ A^0(N)=(\Ind^G_H(N))^G\rightiso N^H,\qquad
	\varphi\mapsto\varphi(1_G). \]
	Untuk memperluas isomorfisme ini secara kanonik menjadi isomorfisme
	funktor-$\delta$ kohomologis $A^n\rightiso\Hm^n(H,\cdot)$, cukup buktikan
	bahwa $A^n$ dapat dihapus untuk $n>0$
	(Definisi~\ref{def:erasable}, Proposisi~\ref{prop:erasable-univ-delta}).
	Menurut Proposisi~\ref{prop:Ind-preservation}, jika $N$ disertakan dalam
	modul-$H$ injektif $I$, maka
	$\Ind^G_H(N)\hookrightarrow\Ind^G_H(I)$ dan $\Ind^G_H(I)$ merupakan
	modul-$G$ injektif. Karena itu $n>0$ mengakibatkan
	$\Hm^n(G,\Ind^G_H(I))=0$, yang membuktikan sifat dapat dihapus.

	Gagasan untuk homologi serupa dan memakai sifat kodapat-dihapus dari
	funktor-$\delta$ homologis. Isomorfisme suku berderajat nol yang diperlukan,
	$N_H\rightiso(\iInd^G_H(N))_G$, juga berasal dari
	Korolari~\ref{prop:Ind-invariant}.

	Metode kedua bekerja dalam kategori turunan. Untuk kohomologi, sekali lagi
	kita mulai dari isomorfisme kanonik
	\[ \underbracket{\Hom_G(\Bbbk,\Ind^G_H(\cdot))}_{
	\simeq\Ind^G_H(\cdot)^G}
	\simeq
	\underbracket{\Hom_H(\Bbbk,\cdot)}_{\simeq(\cdot)^H}
	:H\dcate{Mod}\to\Bbbk\dcate{Mod}. \]
	Kedua ruas merupakan funktor eksak kiri. Karena $\Ind^G_H$ mempertahankan
	objek injektif, dengan mengambil funktor turunan kanan dalam
	Teorema~\ref{prop:localization-triangulated-composite}~(iii), diperoleh
	\begin{equation*}
		\RHom_G(\Bbbk,\mathrm{R}\Ind^G_H(\cdot))
		\simeq\RHom_H(\Bbbk,\cdot).
	\end{equation*}

	Akan tetapi, $\Ind^G_H$ eksak, sehingga $\mathrm{R}\Ind^G_H(\cdot)$
	merupakan funktor yang ditentukan oleh sifat universal lokalisasi dan dapat
	tetap ditulis sebagai $\Ind^G_H$. Jadi, pada tingkat kategori turunan
	terdapat isomorfisme funktor triangulasi
	\[ \RHom_G(\Bbbk,\Ind^G_H(\cdot))
	\simeq\RHom_H(\Bbbk,\cdot). \]
	Mengambil $\Hm^n$ pada kedua ruas memberikan hasil yang diminta.

	Gagasan untuk homologi serupa; unsur kuncinya ialah isomorfisme
	\[ \Bbbk\otimesL[{\Bbbk[G]}]
	\left(\Bbbk[G]\otimesL[{\Bbbk[H]}](\cdot)\right)
	\simeq\Bbbk\otimesL[{\Bbbk[H]}](\cdot). \]
	Ini juga merupakan kasus khusus Contoh~\ref{eg:change-of-rings} atau
	Proposisi~\ref{prop:otimesL-associativity-constraint} mengenai kendala
	asosiatif. Karena $\iInd^G_H$ eksak,
	$\Bbbk[G]\otimesL[{\Bbbk[H]}](\cdot)$ juga dapat diganti dengan
	$\Bbbk[G]\dotimes{\Bbbk[H]}(\cdot)$; langkah lainnya sama.
\end{proof}

% segment-id: o014.li2.chapter6.induced-modules.p005
Perhatikan bahwa jika $\Res^G_H$ ditafsirkan sebagai
$\Bbbk[G]\dotimes{\Bbbk[G]}(\cdot)$ menurut
\eqref{eqn:G-module-pullback-tensor}, maka karena
$\Res^G_H\Bbbk=\Bbbk$, isomorfisme
$\RHom_G(\Bbbk,\mathrm{R}\Ind^G_H(\cdot))
\simeq\RHom_H(\Bbbk,\cdot)$ merupakan relasi adjungsi antara $\RHom$ dan
$\otimesL$ (Teorema~\ref{prop:RHom-otimesL-adjunction}).

% segment-id: o014.li2.chapter6.induced-modules.p006
Kohomologi dan homologi grup dapat dideskripsikan secara konkret melalui
kompleks standar dan kompleks rantai standar dalam
Proposisi~\ref{prop:groupcoh-cochain} dan \ref{prop:grouph-chain}. Kita ingin
membuat isomorfisme dalam Teorema~\ref{prop:Shapiro} eksplisit melalui
kompleks-kompleks tersebut. Tidak ada kesulitan mendasar. Untuk membedakan
grupnya, bagi setiap grup $H$ kita tulis kompleks modul-$H$ $\mathsf{L}$ dari
Definisi~\ref{def:resolution-trivial-module} sebagai $\mathsf{L}^H$; kompleks
ini memberikan resolusi bebas $\mathsf{L}^H\to\Bbbk$.

% segment-id: o014.li2.chapter6.induced-modules.l001
\begin{lemma}\label{prop:Shapiro-alpha-beta}
	Misalkan $\varphi:H\to G$ homomorfisme grup. Homomorfisme terkait
	$H^{n+1}\to G^{n+1}$ menginduksi kuasi-isomorfisme kompleks
	$\mathsf{L}^H\to\varphi^*(\mathsf{L}^G)$, yang komposisinya dengan
	$\varphi^*(\mathsf{L}^G)\to\varphi^*\Bbbk=\Bbbk$ sama dengan
	$\mathsf{L}^H\to\Bbbk$.

	Khususnya, jika $H$ subgrup dari $G$, terdapat kuasi-isomorfisme
	$\mathsf{L}^H\to\Res^G_H\mathsf{L}^G$ dengan sifat komposisi di atas.
\end{lemma}
\begin{proof}
	Dari definisi langsung terlihat bahwa
	$\mathsf{L}^H\to\varphi^*(\mathsf{L}^G)$ merupakan morfisme kompleks;
	pernyataan mengenai komposisinya juga jelas. Karena $\varphi^*$ eksak,
	sedangkan $\mathsf{L}^G\to\Bbbk$ dan $\mathsf{L}^H\to\Bbbk$ keduanya
	kuasi-isomorfisme, maka
	$\mathsf{L}^H\to\varphi^*(\mathsf{L}^G)$ juga kuasi-isomorfisme.
\end{proof}

% segment-id: o014.li2.chapter6.induced-modules.t003
\begin{proposition}\label{prop:Shapiro-std-coh}
	Misalkan $H$ subgrup dari $G$ dan $N$ modul-$H$. Identifikasikan
	$\Hm^n(G,\Ind^G_H(N))$ (atau $\Hm^n(H,N)$) dengan $\Hm^n$ dari kompleks
	standar $C(G,\Ind^G_H(N))$ (atau $C(H,N)$) menurut
	Proposisi~\ref{prop:groupcoh-cochain}. Maka isomorfisme
	\[ \Hm^n\left(G,\Ind^G_H(N)\right)\simeq\Hm^n(H,N),
	\qquad n\in\Z_{\geq0}, \]
	berasal dari morfisme kompleks
	\begin{align*}
		C^m(G,\Ind^G_H(N))&\to C^m(H,N),\\
		[f:G^m\to\Ind^G_H(N)]&\mapsto
		[(h_1,\ldots,h_m)\mapsto f(h_1,\ldots,h_m)(1_G)]\\
		&\quad=(f|_{H^m})(\cdot)(1_G),
	\end{align*}
	dengan $h_1,\ldots,h_m\in H$.
\end{proposition}
\begin{proof}
	Sekali lagi ada dua metode. Pertama, periksalah bahwa pemetaan
	$C^m(G,\Ind^G_H(N))\to C^m(H,N)$ yang ditampilkan memang membentuk
	morfisme kompleks. Karena $\Ind^G_H$ eksak, ketika $N$ bervariasi pemetaan
	ini menginduksi morfisme antara funktor-$\delta$ kohomologis. Pada suku
	berderajat nol, pemetaan tersebut jelas sama dengan isomorfisme kanonik
	$\Ind^G_H(N)^G\rightiso N^H$, $\varphi\mapsto\varphi(1_G)$. Kita telah
	mengetahui bahwa $\Hm^n(G,\Ind^G_H(\cdot))$ merupakan funktor-$\delta$
	kohomologis universal (Teorema~\ref{prop:Shapiro}). Jadi morfisme di atas
	unik dan harus sama dengan isomorfisme
	$\Hm^n(G,\Ind^G_H(N))\simeq\Hm^n(H,N)$ dalam teorema tersebut.

	Untuk metode kedua, misalkan $M$ (berturut-turut $N$) suatu kompleks
	modul-$G$ (berturut-turut modul-$H$) terbatas ke bawah. Perluas funktor
	$\Ind^G_H$ dan $\Res^G_H$ derajat demi derajat ke kompleks, dengan notasi
	yang sama. Relasi adjungsi Proposisi~\ref{prop:pullback-two-adjoint} terangkat
	ke kompleks $\Hom$:
	\[ \Hom^\bullet_H(\Res^G_HM,N)
	\simeq\Hom^\bullet_G(M,\Ind^G_H(N)). \]
	Pemeriksaan langsung tidak sulit; relasi ini juga dapat dipahami sebagai
	adjungsi antara $\Hom^\bullet$ dan $\otimes$ seperti dijelaskan sebelumnya
	(Proposisi~\ref{prop:tensor-Hom-cplx}).

	Sekarang misalkan $N$ modul-$H$ dan ambil resolusi injektif
	$\iota:N\to I$. Maka $\Ind^G_H(N)\to\Ind^G_H(I)$ tetap merupakan resolusi
	injektif, dan naturalitas kompleks $\Hom$ memberikan diagram komutatif
	kompleks berikut (bagian garis penuh):
	\begin{equation}\label{eqn:Shapiro-explicit-diagram}
	\begin{tikzcd}[row sep=small]
		\Hom^\bullet_H(\Bbbk,I) \arrow[d,"\alpha"'] &
		\Hom^\bullet_G(\Bbbk,\Ind^G_H(I)) \arrow[l,"\sim"'] \arrow[d] \\
		\Hom^\bullet_H(\Res^G_H\mathsf{L}^G,I) \arrow[d,"\beta"'] &
		\Hom^\bullet_G(\mathsf{L}^G,\Ind^G_H(I)) \arrow[l,"\sim"] \\
		\Hom^\bullet_H(\mathsf{L}^H,I) & \\
		\Hom^\bullet_H(\mathsf{L}^H,N) \arrow[u] &
		\Hom^\bullet_G(\mathsf{L}^G,\Ind^G_H(N))
		\arrow[uu] \arrow[dashed,l,"\gamma"]
	\end{tikzcd}
	\end{equation}
	dengan $\alpha$ dan $\beta$ berasal dari morfisme dalam
	Lema~\ref{prop:Shapiro-alpha-beta}. Jadi $\beta\alpha$ adalah pemetaan yang
	dikenal $\Hom^\bullet_H(\Bbbk,I)\to\Hom^\bullet_H(\mathsf{L}^H,I)$.
	Isomorfisme horizontal bergaris penuh berasal dari adjungsi, dan semua
	panah vertikal merupakan kuasi-isomorfisme.

	Ingat kembali bukti kedua Teorema~\ref{prop:Shapiro}. Baris pertama
	\eqref{eqn:Shapiro-explicit-diagram} memberikan bentuk dalam kategori
	turunan dari
	\[ \Hm^n(H,N)\simeq\Hm^n(G,\Ind^G_H(N)),\qquad n\in\Z. \]
	Di sisi lain, inti Proposisi~\ref{prop:groupcoh-cochain} adalah isomorfisme
	kompleks
	\[ \Hom^\bullet_H(\mathsf{L}^H,N)\simeq C(H,N),\qquad
	\Hom^\bullet_G(\mathsf{L}^G,\Ind^G_H(N))
	\simeq C(G,\Ind^G_H(N)). \]
	Jadi cukup melengkapi panah $\gamma$ dalam
	\eqref{eqn:Shapiro-explicit-diagram} agar perseginya komutatif dan
	mendeskripsikan panah itu. Karena $\iota$ monik, jika $\gamma$ ada maka ia
	unik.

	Ambil $m\in\Z_{\geq0}$ dan
	$\phi\in\Hom^m_G(\mathsf{L}^G,\Ind^G_H(N))$. Dengan menelaah konstruksinya,
	citra $\phi$ dalam $\Hom^m_H(\mathsf{L}^H,I)$ adalah
	\[ (h_0,\ldots,h_m)\mapsto
	\iota\left(\phi(h_0,\ldots,h_m)(1_G)\right). \]
	Perbandingan dengan isomorfisme dalam bukti
	Proposisi~\ref{prop:groupcoh-cochain} segera menunjukkan bahwa $\gamma$
	dapat didefinisikan seperti yang dinyatakan.
\end{proof}

% segment-id: o014.li2.chapter6.induced-modules.p007
Sekarang kita nyatakan versi homologis. Kita memakai notasi kompleks rantai
standar dari Proposisi~\ref{prop:grouph-chain}.

% segment-id: o014.li2.chapter6.induced-modules.t004
\begin{proposition}\label{prop:Shapiro-std-ho}
	Misalkan $H$ subgrup dari $G$ dan $N$ modul-$H$. Identifikasikan
	$\Hm_n(G,\iInd^G_H(N))$ (atau $\Hm_n(H,N)$) dengan $\Hm_n$ dari kompleks
	rantai standar $C(G,\iInd^G_H(N))$ (atau $C(H,N)$) menurut
	Proposisi~\ref{prop:grouph-chain}. Maka isomorfisme
	\[ \Hm_n\left(G,\iInd^G_H(N)\right)\simeq\Hm_n(H,N),
	\qquad n\in\Z_{\geq0}, \]
	berasal dari morfisme kompleks
	\begin{align*}
		C_m(H,N)&\to C_m(G,\iInd^G_H(N)),\\
		(h_1|\cdots|h_m)^\dagger\otimes y
		&\mapsto(h_1|\cdots|h_m)^\dagger\otimes(1_G\otimes y),
	\end{align*}
	dengan $h_1,\ldots,h_m\in H$ dan $y\in N$.
\end{proposition}
\begin{proof}
	Kita hanya menjelaskan pendekatan kategori turunan. Ambil resolusi projektif
	$\pi:P\to N$ dalam modul-$H$. Diagram komutatif
	\eqref{eqn:Shapiro-explicit-diagram} dari bukti
	Proposisi~\ref{prop:Shapiro-std-coh} harus diganti dengan
	\[\begin{tikzcd}[row sep=small]
		\Bbbk\dotimes{\Bbbk[H]}P \arrow[r,"\sim"] &
		\Bbbk\dotimes{\Bbbk[G]}\iInd^G_H(P) \\
		\Res^G_H\mathsf{L}^G\dotimes{\Bbbk[H]}P
		\arrow[r,"\sim"] \arrow[u,"\alpha"] &
		\mathsf{L}^G\dotimes{\Bbbk[G]}\iInd^G_H(P)
		\arrow[u] \arrow[dd] \\
		\mathsf{L}^H\dotimes{\Bbbk[H]}P \arrow[u,"\beta"] \arrow[d] & \\
		\mathsf{L}^H\dotimes{\Bbbk[H]}N \arrow[dashed,r,"\gamma"'] &
		\mathsf{L}^G\dotimes{\Bbbk[G]}\iInd^G_H(N),
	\end{tikzcd}\]
	dengan panah horizontal bergaris penuh berasal dari kendala asosiatif hasil
	kali tensor pada tingkat kompleks, sedangkan semua panah vertikal merupakan
	kuasi-isomorfisme. Sekali lagi, tinggal mendeskripsikan $\gamma$.
	Langkah selanjutnya serupa dengan versi kohomologis.
\end{proof}
